

\chapter{Introduction Générale}

\section{Présentation de l’organisme d’accueil}
\subsection{Faculté de Médecine de Monastir (FMM)}
Fondée en 1980, la Faculté de Médecine de Monastir (FMM) est un établissement public d’enseignement supérieur et de recherche médicale en Tunisie. Elle a pour mission de former des professionnels de santé compétents et de promouvoir la recherche médicale appliquée.

\textbf{Certifications et collaborations} :
\begin{itemize}
  \item ISO 21001:2018 pour la gestion de la qualité de l’éducation.
  \item Partenariats internationaux avec des hôpitaux universitaires européens.
\end{itemize}

\subsection{SmartLab}
Le SmartLab est un laboratoire d’innovation spécialisé dans le développement de solutions médicales low-cost. Il met à disposition :
\begin{itemize}
  \item Des dispositifs embarqués (Raspberry Pi, capteurs biométriques).
  \item Des serveurs de traitement des données sécurisés.
\end{itemize}

\section{Contexte et justification du projet}
\subsection{Problématique}
Dans les salles de réveil tunisiennes, la surveillance des patients est majoritairement manuelle. Les infirmiers doivent vérifier régulièrement les constantes vitales sur des machines anciennes non connectées. Cela entraîne :
\begin{itemize}
  \item Des retards d’intervention supérieurs à 5 minutes.
  \item Un taux d’erreurs humaines de plus de 15\%.
\end{itemize}

\subsection{Opportunité}
Face au coût élevé des solutions commerciales connectées, une approche par vision artificielle est envisagée. Celle-ci permet de :
\begin{itemize}
  \item Réutiliser les machines existantes.
  \item Numériser les informations via OCR.
\end{itemize}

\section{Objectifs}
\subsection{Objectif principal}
Développer une solution low-cost de surveillance automatisée en salle de réveil, basée sur la capture d’image et la reconnaissance optique de caractères (OCR).

\subsection{Objectifs spécifiques}
\begin{itemize}
  \item Extraire les constantes vitales (SpO2, FC) à l’aide d’OpenCV et Tesseract.
  \item Générer des alertes via Firebase Cloud Messaging et SMS.
  \item Réduire la latence d’alerte à moins de 3 secondes.
  \item Maintenir un budget inférieur à 10 000€.
\end{itemize}

\chapter{Cadre de Projet}

\section{Portée et livrables}
\subsection{Périmètre}
\begin{itemize}
  \item \textbf{Inclus} : Surveillance de 10 lits avec caméras IR.
  \item \textbf{Exclus} : Intégration à des moniteurs connectés.
\end{itemize}

\subsection{Livrables}
\begin{itemize}
  \item Caméras Raspberry Pi 5 avec module infrarouge.
  \item Backend Python (Flask) + base SQLite.
  \item Interface web (React.js).
\end{itemize}

\section{Analyse de l’existant}
\subsection{Solutions alternatives}
\begin{tabular}{|l|l|l|}
\hline
\textbf{Méthode} & \textbf{Avantages} & \textbf{Inconvénients} \\
\hline
Caméra + OCR & Économique, sans contact & Sensible à l’éclairage \\
\hline
Capteurs DIY & Haute précision & Coût, complexité technique \\
\hline
\end{tabular}

\subsection{Benchmark OCR}
\begin{itemize}
  \item Tesseract 5.0 avec prétraitement (binarisation, filtrage).
  \item Taux de reconnaissance : 93\% sur écrans LCD.
\end{itemize}

\section{Risques et atténuation}
\begin{itemize}
  \item \textbf{Qualité d’image} : Solution -- caméra IR et filtrage logiciel (OpenCV).
  \item \textbf{Confidentialité} : Chiffrement AES-256 et stockage local.
\end{itemize}

\chapter{Méthodologie de Travail}

\section{Méthodologie Agile adaptée}
\subsection{Cycle de développement}
\begin{itemize}
  \item Sprint 0 : Cadrage du besoin (2 semaines).
  \item Sprint 1 : Extraction OCR (SpO2, FC).
  \item Sprint 2 : Mécanisme d’alerte (FCM + Twilio).
\end{itemize}

\subsection{Rôles}
\begin{itemize}
  \item \textbf{Product Owner} : Médecin coordinateur.
  \item \textbf{Développeurs} :
  \begin{itemize}
    \item 1 ingénieur Python.
    \item 1 développeur React.
  \end{itemize}
\end{itemize}

\section{Outils utilisés}
\begin{itemize}
  \item Trello pour la gestion de projet (Kanban).
  \item GitLab pour le versionnage et la CI/CD.
\end{itemize}

\chapter{Étude Architecturale}

\section{Architecture logicielle}
\subsection{Flux de traitement}
\begin{enumerate}
  \item Capture image via caméra.
  \item Détection du cadre de l’écran (OpenCV).
  \item Extraction texte (Tesseract).
  \item Déclenchement d’alerte si dépassement de seuils.
\end{enumerate}

\subsection{Technologies utilisées}
\begin{itemize}
  \item \textbf{Backend} : Flask, PyTesseract, Firebase SDK.
  \item \textbf{Frontend} : React.js, Material-UI.
\end{itemize}

\section{Sécurité des données}
\begin{itemize}
  \item Chiffrement AES-256 pour les données.
  \item Authentification via JWT.
\end{itemize}

\section{Prototype matériel}
\begin{itemize}
  \item Caméras 1080p avec support ajustable.
  \item Installation discrète au-dessus des écrans.
\end{itemize}

\appendix

\chapter*{Annexes}
\addcontentsline{toc}{chapter}{Annexes}
\begin{itemize}
  \item \textbf{Annexe A} : Extrait OCR (SpO2 = 98, FC = 72).
  \item \textbf{Annexe B} : Schéma de la base de données SQLite.
  \item \textbf{Annexe C} : Guide d’installation des caméras.
\end{itemize}

% \end{document}
